\section{Variance Reduction Techniques}

\subsection{Control Variates}

\subsubsection*{Basic method}

Let $Y_1,\ldots,Y_n$ be i.i.d. outputs for estimating $\mathbb{E}[Y]$, and suppose $X_i$ is generated with $Y_i$ and $\mathbb{E}[X]$ is known. For fixed $b$ define
$Y_i(b)=Y_i-b(X_i-\mathbb{E}[X])$,

and estimate
$\bar{Y}(b)=\bar{Y}-b(\bar{X}-\mathbb{E}[X])=\frac{1}{n}\sum_{i=1}^n\{Y_i-b(X_i-\mathbb{E}[X])\}$.

This is unbiased for all $b$. Its per-replication variance is

$\sigma^2(b)=\operatorname{var}[Y_i(b)]=\sigma_Y^2-2b\sigma_X\sigma_Y\rho_{XY}+b^2\sigma_X^2$.

The optimal coefficient is
$b^*=\frac{\sigma_Y}{\sigma_X}\rho_{XY}=\frac{\operatorname{Cov}(X,Y)}{\operatorname{var}(X)}$.

The resulting variance ratio relative to plain Monte Carlo is

$\frac{\operatorname{var}[\bar{Y}-b^*(\bar{X}-\mathbb{E}[X])]}{\operatorname{var}[\bar{Y}]}=1-\rho_{XY}^2$.


\subsubsection*{Estimated coefficient and CI}

In practice use the sample regression coefficient
$\hat{b}_n=\frac{\sum_{i=1}^n(X_i-\bar{X})(Y_i-\bar{Y})}{\sum_{i=1}^n(X_i-\bar{X})^2}$.

Then $\hat{b}_n\to b^*$ a.s. and
$\bar{Y}(\hat{b}_n)=\frac{1}{n}\sum_{i=1}^n\{Y_i-\hat{b}_n(X_i-\mathbb{E}[X])\}$.

Using $\hat{b}_n$ creates finite-sample bias, but has the same asymptotic precision as using $b^*$. For fixed $b$,
$\frac{\bar{Y}(b)-\mathbb{E}[Y]}{\sigma(b)/\sqrt{n}}\Rightarrow N(0,1)$,
with sample standard deviation

$s(b)=\left[\frac{1}{n-1}\sum_{i=1}^n(Y_i(b)-\bar{Y}(b))^2\right]^{1/2}$.

Hence an asymptotic CI is $\bar{Y}(b)\pm z_{\delta/2}s(b)/\sqrt{n}$. Also
$\frac{\bar{Y}(\hat{b}_n)-\mathbb{E}[Y]}{s(\hat{b}_n)/\sqrt{n}}\Rightarrow N(0,1)$.

\subsubsection*{Financial controls}

Any martingale with known initial value is a natural control. Under risk-neutral pricing with constant $r$, $e^{-rt}S(t)$ is a martingale, so $\mathbb{E}[e^{-rT}S(T)]=S(0)$. For discounted payoff $Y$,
$\frac{1}{n}\sum_{i=1}^n\{Y_i-b[e^{-rT}S_i(T)-S(0)]\}$

uses the stock as control. For a call, this works best near $K=0$ and may weaken for deep OTM strikes. A stronger control can be a tractable option. For Asian options,
$\bar{S}_A=\frac{1}{n}\sum_{i=1}^nS(t_i),\qquad \bar{S}_G=\left(\prod_{i=1}^nS(t_i)\right)^{1/n}$.

Under GBM, $\bar{S}_G$ is lognormal and the geometric-average Asian option can be priced analytically; its payoff is often highly correlated with the arithmetic-average Asian payoff.

\subsection{Antithetic Variates}

\subsubsection*{Estimator}

Antithetic variates pair inputs to create negative dependence. If $U\sim\operatorname{Unif}[0,1]$, use $U$ and $1-U$; for a symmetric normal input, use $Z$ and $-Z$.

Let $(Y_i,\tilde{Y}_i)$ be i.i.d. pairs with the same marginal distribution. The estimator is

$\hat{Y}_{AV}=\frac{1}{2n}\left(\sum_{i=1}^nY_i+\sum_{i=1}^n\tilde{Y}_i\right)=\frac{1}{n}\sum_{i=1}^n\frac{Y_i+\tilde{Y}_i}{2}$.

Apply the CLT to pair averages:
$\frac{\hat{Y}_{AV}-\mathbb{E}[Y]}{\sigma_{AV}/\sqrt{n}}\Rightarrow N(0,1),\quad \sigma_{AV}^2=\operatorname{var}\left(\frac{Y_i+\tilde{Y}_i}{2}\right)$.

A CI is $\hat{Y}_{AV}\pm z_{\delta/2}s_{AV}/\sqrt{n}$, where $s_{AV}$ is the sample sd of pair averages.

Compare $n$ antithetic pairs with $2n$ independent samples. Since

$\operatorname{var}[Y_i+\tilde{Y}_i]=2\operatorname{var}[Y_i]+2\operatorname{Cov}(Y_i,\tilde{Y}_i)$,

antithetic sampling reduces variance iff $\operatorname{Cov}(Y_i,\tilde{Y}_i)<0$. Monotone outputs often satisfy this when paired by $U,1-U$ or $Z,-Z$.

\subsubsection*{GBM call example}

For a call, pair $Z_i$ and $-Z_i$:

$S_i^+=S_0\exp\left((r-\frac12\sigma^2)T+\sigma\sqrt{T}Z_i\right)$,

$S_i^-=S_0\exp\left((r-\frac12\sigma^2)T-\sigma\sqrt{T}Z_i\right)$.

Set $X_i=e^{-rT}(S_i^+-K)^+$ and $Y_i=e^{-rT}(S_i^--K)^+$. Then
$\hat{v}=\frac{1}{n}\sum_{i=1}^n\frac{X_i+Y_i}{2}$,

$\mathrm{SE}=\left[\frac{1}{n(n-1)}\left\{\sum_{i=1}^n\left(\frac{X_i+Y_i}{2}\right)^2-n\hat{v}^2\right\}\right]^{1/2}$.

\subsection{Stratified Sampling}

\subsubsection*{Basic method}

Let $A_1,\ldots,A_K$ be disjoint strata with probabilities $p_i$. Then

$\mathbb{E}[Y]=\sum_{i=1}^Kp_i\mathbb{E}[Y\mid Y\in A_i]$.

With proportional allocation $n_i=np_i$ and conditional draws $Y_{ij}\sim (Y\mid Y\in A_i)$,

$\hat{Y}=\sum_{i=1}^Kp_i\frac{1}{n_i}\sum_{j=1}^{n_i}Y_{ij}=\frac{1}{n}\sum_{i=1}^K\sum_{j=1}^{n_i}Y_{ij}$.

\subsubsection*{General allocation}

Often strata are defined by another variable $X$. If $A_i$ partition the support of $X$,

$\mathbb{E}[Y]=\sum_{i=1}^Kp_i\mathbb{E}[Y\mid X\in A_i]$.

For $n_i\ge1$, $\sum_i n_i=n$, and $q_i=n_i/n$,

$\hat{Y}=\sum_{i=1}^Kp_i\frac{1}{n_i}\sum_{j=1}^{n_i}Y_{ij}=\frac{1}{n}\sum_{i=1}^K\frac{p_i}{q_i}\sum_{j=1}^{n_i}Y_{ij}$.

Let $\mu_i=\mathbb{E}[Y\mid X\in A_i]$ and $\sigma_i^2=\operatorname{var}(Y\mid X\in A_i)$. Then $\hat{Y}$ is unbiased and

$\operatorname{var}[\hat{Y}]=\sum_{i=1}^Kp_i^2\frac{\sigma_i^2}{n_i}=\frac{\sigma^2(q)}{n},\quad \sigma^2(q)=\sum_{i=1}^K\frac{p_i^2}{q_i}\sigma_i^2$.
With fixed $K$, $\sqrt{n}(\hat{Y}-\mu)\Rightarrow N(0,\sigma^2(q))$, and
$s^2(q)=\sum_{i=1}^K\frac{p_i^2}{q_i}s_i^2$
estimates $\sigma^2(q)$.

\subsubsection*{Uniform and inverse-transform strata}

For $U\sim\operatorname{Unif}[0,1]$, equal strata are $A_i=((i-1)/n,i/n)$. With $U_i\sim\operatorname{Unif}[0,1]$,

$V_i=\frac{i-1}{n}+\frac{U_i}{n},\qquad \hat{Y}=\frac{1}{n}\sum_{i=1}^nf(V_i)$.

For a distribution with CDF $F$, choose probabilities $p_i$ and quantiles $a_i=F^{-1}(p_1+\cdots+p_i)$. Conditional simulation in stratum $(a_{i-1},a_i]$ uses

$V=F(a_{i-1})+U\{F(a_i)-F(a_{i-1})\},\qquad F^{-1}(V)\sim (Y\mid Y\in A_i)$.

\subsubsection*{Proportional and optimal allocation}

For proportional allocation $q_i=p_i$,
$\sigma^2(q)=\sum_{i=1}^Kp_i\sigma_i^2$.

Plain variance decomposes as

$\operatorname{var}(Y)=\sum_{i=1}^Kp_i\sigma_i^2+\sum_{i=1}^Kp_i\mu_i^2-\left(\sum_{i=1}^Kp_i\mu_i\right)^2$.

By Jensen, proportional stratification cannot increase variance, with strict gain unless all $\mu_i$ are equal. Optimizing $\sigma^2(q)$ gives

$q_i^*=\frac{p_i\sigma_i}{\sum_{\ell=1}^Kp_\ell\sigma_\ell},\quad \sigma^2(q^*)=\left(\sum_{i=1}^Kp_i\sigma_i\right)^2$.


\subsubsection*{Call option strata}

For a GBM call, write $W_T=\sqrt{T}\Phi^{-1}(U)$ with $U\sim\operatorname{Unif}[0,1]$ and

$h(u)=\left(S_0\exp\left\{-\frac12\sigma^2T+\sigma\sqrt{T}\Phi^{-1}(u)\right\}-e^{-rT}K\right)^+$.

With $k$ equal strata, $U_{ij}=(i-1)/k+V_{ij}/k$, $\hat{\mu}_i=n_i^{-1}\sum_{j=1}^{n_i}h(U_{ij})$, and

$\hat{\mu}=\frac{1}{k}\sum_{i=1}^k\hat{\mu}_i,\qquad \mathrm{SE}=\frac{1}{k}\left(\sum_{i=1}^k\frac{s_i^2}{n_i}\right)^{1/2}$.

\subsection{Importance Sampling}

\subsubsection*{Change of measure}

To estimate $\mu=\mathbb{E}[h(X)]$ with density $f$, choose an alternative density $g$ such that $g(x)=0$ implies $f(x)=0$. If $Y\sim g$,

$\mu=\int h(x)f(x)\,dx=\int h(x)\frac{f(x)}{g(x)}g(x)\,dx=\mathbb{E}_g\left[h(Y)\frac{f(Y)}{g(Y)}\right]$.

The likelihood ratio is $f(Y)/g(Y)$, and the estimator averages $h(Y_i)f(Y_i)/g(Y_i)$. In a discrete model, use $h(Y_i)p(Y_i)/\bar{p}(Y_i)$.

\subsubsection*{Variance and zero-variance density}

The IS estimator has variance
$\frac{1}{n}\operatorname{var}_g\left[h(Y)\frac{f(Y)}{g(Y)}\right]$.

For $h\ge0$, the zero-variance density is

$g^*(x)=c\,h(x)f(x),\qquad \frac{1}{c}=\int h(x)f(x)\,dx=\mu$.

It is unusable directly because it requires $\mu$, but it says $g$ should mimic $h(x)f(x)$.

\subsubsection*{Normal mean shift and options}

For $X\sim N(0,1)$, choose $Y\sim N(\theta,1)$. The likelihood ratio is

$\frac{f(x)}{g(x)}=\exp\left(-\theta x+\frac12\theta^2\right)$.

Mode matching sets $\theta=\arg\max_x h(x)f(x)$. For a binary call,

$\mathbb{E}[e^{-rT}\mathbf{1}_{\{S_T\ge K\}}]=\mathbb{E}[e^{-rT}\mathbf{1}_{\{X\ge b\}}]$,
where
$b=\frac{\log(K/S_0)-(r-\sigma^2/2)T}{\sigma\sqrt{T}}$.

Here $h(x)=e^{-rT}\mathbf{1}_{\{x\ge b\}}$, $x^*=\max\{b,0\}$, and with $Y\sim N(x^*,1)$ the output is

$h(Y)\frac{f(Y)}{g(Y)}=\begin{cases}0,&Y<b,\\ \exp\{-rT-x^*Y+(x^*)^2/2\},&Y\ge b.\end{cases}$

For a vanilla call,
$h(x)=\left(S_0\exp\left\{-\frac12\sigma^2T+\sigma\sqrt{T}x\right\}-e^{-rT}K\right)^+$,

and the shift pushes sampling toward exercise states.

\subsubsection*{Exponential tilting}

For density $f$, the exponential tilt family is

$f_\theta(x)=\frac{e^{\theta x}f(x)}{\mathbb{E}[e^{\theta X}]},\, \theta\in\mathbb{R}$.
For a pmf $p$,
$p_\theta(x)=\frac{e^{\theta x}p(x)}{\mathbb{E}[e^{\theta X}]}$.

For normal $X$, this is mean shifting. For rare large losses, take $\theta>0$.

\subsubsection*{Bernoulli rare event}

Let independent defaults satisfy $P(X_k=1)=p_k$ and $L=\sum_{k=1}^mX_k$. To estimate $P(L>x)$ with $x>\mathbb{E}[L]$, tilt each Bernoulli:

$\bar{p}_k=P(Y_k=1)=\frac{p_ke^\theta}{1+p_k(e^\theta-1)},\qquad \bar{L}=\sum_{k=1}^mY_k$.

The IS output is
$H=\mathbf{1}_{\{\bar{L}>x\}}\prod_{k=1}^m\left(\frac{p_k}{\bar{p}_k}\right)^{Y_k}\left(\frac{1-p_k}{1-\bar{p}_k}\right)^{1-Y_k}$.

With
$\phi(\theta)=\sum_{k=1}^m\log[1+p_k(e^\theta-1)]$,

one has $E[H^2]\le e^{-2\theta x+2\phi(\theta)}$. Choose $\theta>0$ by solving $\phi'(\theta)=x$.

\subsection{Cross-Entropy Method}

\subsubsection*{Principle}

Cross-entropy selects a parametric IS density $f_\theta$ close to the zero-variance density
$g^*(x)=\frac{1}{\mu}h(x)f(x)$.

Minimize
$R(g^*\|f_\theta)=\int \log\left(\frac{g^*(x)}{f_\theta(x)}\right)g^*(x)\,dx$,

equivalently maximize
$\int h(x)f(x)\log f_\theta(x)\,dx$.

The population first-order equation is
$\mathbb{E}\left[h(X)\frac{\partial}{\partial\theta}\log f_\theta(X)\right]=0$.

\subsubsection*{Basic algorithm}

Using pilot samples $X_1,\ldots,X_N\sim f$, solve
$\frac{1}{N}\sum_{k=1}^Nh(X_k)\frac{\partial}{\partial\theta}\log f_\theta(X_k)=0$.

For $f=N(0,I_m)$ and $f_\theta=N(\theta,I_m)$,
$\hat{\theta}=\frac{\sum_{k=1}^Nh(X_k)X_k}{\sum_{k=1}^Nh(X_k)}$.

Then sample $Y_i\sim f_{\hat{\theta}}$ and average
$H_i=h(Y_i)\frac{f(Y_i)}{f_{\hat{\theta}}(Y_i)}$.

For a one-dimensional normal shift, $H_i=h(Y_i)\exp\{-\hat{\theta}Y_i+\hat{\theta}^2/2\}$. For very rare events, the basic method may fail because $\sum_kh(X_k)=0$ in the pilot run.

\subsubsection*{Iterative cross-entropy}

Let $\ell_\theta(x)=f(x)/f_\theta(x)$. If pilot samples are drawn from $f_\nu$, then maximize

$\mathbb{E}_\nu[h(Y)\ell_\nu(Y)\log f_\theta(Y)]$,
and solve
$\mathbb{E}_\nu\left[h(Y)\ell_\nu(Y)\frac{\partial}{\partial\theta}\log f_\theta(Y)\right]=0$.

At iteration $j$, sample $Y_k\sim f_{\hat{\theta}^j}$ and update from

$\frac{1}{N}\sum_{k=1}^Nh(Y_k)\ell_{\hat{\theta}^j}(Y_k)\frac{\partial}{\partial\theta}\log f_\theta(Y_k)=0$.

For $N(\theta,I_m)$,
$\hat{\theta}^{j+1}=\frac{\sum_{k=1}^Nh(Y_k)\ell_{\hat{\theta}^j}(Y_k)Y_k}{\sum_{k=1}^Nh(Y_k)\ell_{\hat{\theta}^j}(Y_k)}$.

Here
$\ell_{\hat{\theta}^j}(Y_k)=\exp(-\langle\hat{\theta}^j,Y_k\rangle+\|\hat{\theta}^j\|^2/2)$.
